% !TeX root = ../main.tex
\begin{proposalbody}
\noindent\textbf{申请学位实践成果}开题报告应包括但不限于：

\ProposalHeading{1}{预期成果背景与意义、成果形式}
% 在这里填写正文；段落之间留一个空行。
\WritingSpace

\ProposalHeading{2}{技术创新性与可行性}
\WritingSpace

\ProposalHeading{3}{市场需求与应用前景}
\WritingSpace

\ProposalHeading{4}{实施路径与风险控制}
\WritingSpace

\ProposalHeading{5}{社会价值与可持续性}
\WritingSpace

\ProposalHeading{6}{研究进度安排}
\begin{proposalschedule}
  \ScheduleRow{}{}{}
  \ScheduleRow{}{}{}
  \ScheduleRow{}{}{}
  \ScheduleRow{}{}{}
  \ScheduleRow{}{}{}
  \ScheduleRow{}{}{}
\end{proposalschedule}

\ProposalHeading{7}{与选题有关的工作积累、已有的研究工作成绩}
\ifTemplateBlanks\vspace{7mm}\fi
\end{proposalbody}
\clearpage
